\section{Conclusiones y trabajo futuro}
\label{sec:conclusiones}

El sistema final responde 37 de 38 preguntas respondibles, se abstiene en las 5 sin respuesta y acierta 33 de 37 en revisi\'on manual. En conversaci\'on alcanza 14 de 14 con compuerta de dominio y 11 de 14 sin ella. RAGAS reporta 0,952 de fidelidad, 0,881 de precisi\'on del contexto y 0,921 de cobertura. La latencia media es de 6,8 s sin razonamiento y 6,5 s con razonamiento, con un modelo local de 7,5B y un reranker en CPU.

La primera lecci\'on es metodol\'ogica: el piloto sobreestima la recuperaci\'on. Con el mismo codificador y las mismas preguntas, Recall@10 cae de 0,947 a 0,684 al pasar de 1.932 a 161.868 documentos. Cualquier comparaci\'on de codificadores o de chunking hecha solo sobre una muestra peque\~na habr\'ia elegido la configuraci\'on equivocada. La segunda es que el reordenamiento rinde mejor sobre un conjunto corto: reordenar el top-5 sube el nDCG 5,6 puntos y el P@1 13,1 sin perder cobertura, mientras que reordenar 20 o 50 candidatos pierde recall y multiplica la latencia. La tercera es que la limpieza conservadora basta: ninguna variante supera a otra de forma significativa, y la ganancia real es eliminar p\'aginas de error, no reescribir el texto.

La cuarta lecci\'on es que la abstenci\'on barata es la compuerta. Rechazar antes de generar cuesta decenas de milisegundos y evita que el modelo improvise sobre temas ajenos. La quinta es que la evaluaci\'on autom\'atica necesita cobertura: una media de RAGAS con 4 valores de 38 invirti\'o una conclusi\'on sobre la precisi\'on del contexto. Reportar $n$ junto a cada media no es un detalle de estilo; cambia lo que se concluye.

El trabajo futuro tiene cuatro frentes. Primero, cerrar los dos fallos de recuperaci\'on: g-046 exige revisar la consulta o el \'indice, y g-031 requiere un conjunto de candidatos mayor con un coste que hoy no se justifica. Segundo, mejorar la correcci\'on de la respuesta, la m\'etrica m\'as d\'ebil de RAGAS. Tercero, completar la revisi\'on del gold con la cuarta anotadora y ampliar el conjunto de prueba para reducir los intervalos de confianza. Cuarto, probar un chunking distinto y una reindexaci\'on completa, que es la \'unica v\'ia que queda para atacar la p\'erdida de contexto sin sacrificar cobertura.
